\documentclass{amsart}
% For dealing with references we use the comment environment
\usepackage{verbatim}
\newenvironment{reference}{\comment}{\endcomment}
%\newenvironment{reference}{}{}
\newenvironment{slogan}{\comment}{\endcomment}
\newenvironment{history}{\comment}{\endcomment}

% For commutative diagrams we use Xy-pic
\usepackage[all]{xy}

% We use 2cell for 2-commutative diagrams.
\xyoption{2cell}
\UseAllTwocells

% We use multicol for the list of chapters between chapters
\usepackage{multicol}

% This is generally recommended for better output
\usepackage{lmodern}
\usepackage[T1]{fontenc}

% For cross-file-references

% Package for hypertext links:
\usepackage{hyperref}

% For any local file, say "hello.tex" you want to link to please

% Theorem environments.
%
\theoremstyle{plain}
\newtheorem{theorem}[subsection]{Theorem}
\newtheorem{proposition}[subsection]{Proposition}
\newtheorem{lemma}[subsection]{Lemma}

\theoremstyle{definition}
\newtheorem{definition}[subsection]{Definition}
\newtheorem{example}[subsection]{Example}
\newtheorem{exercise}[subsection]{Exercise}
\newtheorem{situation}[subsection]{Situation}

\theoremstyle{remark}
\newtheorem{remark}[subsection]{Remark}
\newtheorem{remarks}[subsection]{Remarks}

\numberwithin{equation}{subsection}

% Macros
%
\def\lim{\mathop{\mathrm{lim}}\nolimits}
\def\colim{\mathop{\mathrm{colim}}\nolimits}
\def\Spec{\mathop{\mathrm{Spec}}}
\def\Hom{\mathop{\mathrm{Hom}}\nolimits}
\def\Ext{\mathop{\mathrm{Ext}}\nolimits}
\def\SheafHom{\mathop{\mathcal{H}\!\mathit{om}}\nolimits}
\def\SheafExt{\mathop{\mathcal{E}\!\mathit{xt}}\nolimits}
\def\Sch{\mathit{Sch}}
\def\Mor{\mathop{\mathrm{Mor}}\nolimits}
\def\Ob{\mathop{\mathrm{Ob}}\nolimits}
\def\Sh{\mathop{\mathit{Sh}}\nolimits}
\def\NL{\mathop{N\!L}\nolimits}
\def\CH{\mathop{\mathrm{CH}}\nolimits}
\def\proetale{{pro\text{-}\acute{e}tale}}
\def\etale{{\acute{e}tale}}
\def\QCoh{\mathit{QCoh}}
\def\Ker{\mathop{\mathrm{Ker}}}
\def\Im{\mathop{\mathrm{Im}}}
\def\Coker{\mathop{\mathrm{Coker}}}
\def\Coim{\mathop{\mathrm{Coim}}}

% Boxtimes
%
\DeclareMathSymbol{\boxtimes}{\mathbin}{AMSa}{"02}

%
% Macros for moduli stacks/spaces
%
\def\QCohstack{\mathcal{QC}\!\mathit{oh}}
\def\Cohstack{\mathcal{C}\!\mathit{oh}}
\def\Spacesstack{\mathcal{S}\!\mathit{paces}}
\def\Quotfunctor{\mathrm{Quot}}
\def\Hilbfunctor{\mathrm{Hilb}}
\def\Curvesstack{\mathcal{C}\!\mathit{urves}}
\def\Polarizedstack{\mathcal{P}\!\mathit{olarized}}
\def\Complexesstack{\mathcal{C}\!\mathit{omplexes}}
% \Pic is the operator that assigns to X its picard group, usage \Pic(X)
% \Picardstack_{X/B} denotes the Picard stack of X over B
% \Picardfunctor_{X/B} denotes the Picard functor of X over B
\def\Pic{\mathop{\mathrm{Pic}}\nolimits}
\def\Picardstack{\mathcal{P}\!\mathit{ic}}
\def\Picardfunctor{\mathrm{Pic}}
\def\Deformationcategory{\mathcal{D}\!\mathit{ef}}

\setlength{\emergencystretch}{3em}
\begin{document}
\title[Stacks Project, Tag 02W8]{489 --- Separated locally quasi-finite scheme morphisms satisfy effective fppf descent}
\maketitle
\noindent Copyright (C) 2005--2025 Johan de Jong. Stacks Project authors. Source: \href{https://stacks.math.columbia.edu/tag/02W8}{Tag 02W8}.

\noindent Editorial difficulty note (not author text). Difficulty H2: A locally quasi-finite object need not itself be quasi-compact, so quasi-affine descent cannot be applied globally. The proof saturates each affine open under the descent isomorphism, proves the saturation quasi-compact and invariant using the cocycle, descends these quasi-affine pieces and glues their effective descent data..

\noindent Editorial provenance note (not author text). History: excerpt from revision \texttt{a0444\allowbreak 6e57e\allowbreak c1fbc\allowbreak 252a8\allowbreak 71afc\allowbreak ec775\allowbreak 2fb28\allowbreak 07b14}, more-morphisms.tex, lines 16278--16356. Reference presentation was adapted; original-author context and core support blocks were retained or added. The mathematical wording is unchanged. Licensed under GNU FDL 1.2 or later; no invariant sections or cover texts. See ../licenses/COPYING. Context blocks reproduce author definitions and supporting results; they are not additional counted problems.

\begin{definition}
\label{topologies-definition-fppf-covering}
Let $T$ be a scheme. An {\it fppf covering of $T$} is a family
of morphisms $\{f_i : T_i \to T\}_{i \in I}$ of schemes
such that each $f_i$ is flat, locally of finite presentation and such
that $T = \bigcup f_i(T_i)$.
\end{definition}

\begin{lemma}
\label{lemma-separated-locally-quasi-finite-morphisms-fppf-descend}
Let $S$ be a scheme.
Let $\{X_i \to S\}_{i\in I}$ be an fppf covering, see
Topologies, Definition \ref{topologies-definition-fppf-covering}.
Let $(V_i/X_i, \varphi_{ij})$ be a descent datum
relative to $\{X_i \to S\}$. If each morphism
$V_i \to X_i$ is separated and locally quasi-finite,
then the descent datum is effective.
\end{lemma}

\begin{proof}
Being separated and being locally quasi-finite
are properties of morphisms of schemes
which are preserved under any base change, see
Schemes, Lemma \href{https://stacks.math.columbia.edu/tag/01KU}{lemma separated permanence (Tag 01KU)} and
Morphisms, Lemma \href{https://stacks.math.columbia.edu/tag/01TM}{lemma base change quasi finite (Tag 01TM)}.
Hence Descent, Lemma \href{https://stacks.math.columbia.edu/tag/02W3}{lemma descending types morphisms (Tag 02W3)}
applies and it suffices to prove the statement of the lemma
in case the fppf-covering is given by a single
$\{X \to S\}$ flat surjective morphism of finite presentation of affines.
Say $X = \Spec(A)$ and $S = \Spec(R)$ so
that $R \to A$ is a faithfully flat ring map.
Let $(V, \varphi)$ be a descent datum relative to $X$ over $S$
and assume that $\pi : V \to X$ is separated and
locally quasi-finite.

\medskip\noindent
Let $W^1 \subset V$ be any affine open.
Consider $W = \text{pr}_1(\varphi(W^1 \times_S X)) \subset V$.
Here is a picture
$$
\xymatrix{
W^1 \times_S X \ar[rrrrr] \ar[ddd] \ar[rd]
& & & & &
\varphi(W^1 \times_S X) \ar[ddd] \ar[ld] \\
& V \times_S X \ar[rrr]^\varphi \ar[rd] \ar[dd]
& & &
X \times_S V \ar[ld] \ar[dd] & \\
& &
X \times_S X \ar[r]^1 \ar[d]_{\text{pr}_0}
&
X \times_S X \ar[d]^{\text{pr}_1}
& & \\
W^1 \ar[r] &
V \ar[r] &
X &
X &
V \ar[l] &
W \ar[l]
}
$$
Ok, and now since $X \to S$ is flat and of finite presentation it
is universally open (Morphisms, Lemma \href{https://stacks.math.columbia.edu/tag/01UA}{lemma fppf open (Tag 01UA)}).
Hence we conclude that $W$ is open. Moreover, it is
also clearly the case that $W$ is quasi-compact, and
$W^1 \subset W$. Moreover, we note that
$\varphi(W \times_S X) = X \times_S W$ by the cocycle
condition for $\varphi$. Hence we obtain a new descent datum
$(W, \varphi')$ by restricting $\varphi$ to $W \times_S X$.
Note that the morphism $W \to X$ is quasi-compact, separated
and locally quasi-finite. This implies that it is
separated and quasi-finite by definition. Hence it is quasi-affine by
Lemma \href{https://stacks.math.columbia.edu/tag/02LR}{lemma quasi finite separated quasi affine (Tag 02LR)}.
Thus by
Descent, Lemma \href{https://stacks.math.columbia.edu/tag/0247}{lemma quasi affine (Tag 0247)}
we see that the descent datum
$(W, \varphi')$ is effective.

\medskip\noindent
In other words, we find that there exists an open covering
$V = \bigcup W_i$ by quasi-compact opens $W_i$ which are
stable for the descent morphism $\varphi$.
Moreover, for each such quasi-compact open $W \subset V$
the corresponding descent data $(W, \varphi')$ is effective.
This means the original descent datum is effective by glueing the
schemes obtained from descending the opens $W_i$, see
Descent, Lemma \href{https://stacks.math.columbia.edu/tag/0AP4}{lemma effective for fpqc is local upstairs (Tag 0AP4)}.
\end{proof}
\end{document}
